% ----------------------------------------------------------
\chapter{Diagrama}
\label{cha:diagrama}
\index{Diagrama}
% ----------------------------------------------------------
\begin{figure}[H]
    \centering
		    \caption{Diagrama de arquitetura em camadas (Java + Python)}
    \label{fig:meu_diagrama}
    \begin{tikzpicture}[
        layer/.style={draw, rounded corners, minimum width=8cm, 
                      minimum height=1.2cm, align=center, fill=#1},
        arrow/.style={-{Stealth}, thick}
    ]
        \node[layer=blue!15] (apresentacao) {Camada de Apresentação\\(Interface / API)};
        \node[layer=green!15, below=0.8cm of apresentacao] (negocio) {Camada de Lógica\\(Java + Python)};
        \node[layer=orange!15, below=0.8cm of negocio] (dados) {Camada de Dados\\(CSV / TXT)};
        
        \draw[arrow] (apresentacao) -- (negocio);
        \draw[arrow] (negocio) -- (dados);
    \end{tikzpicture}
	\caption*{Fonte: \oautor}
	\label{fig:diagramaarquitetura}
\end{figure}


%---------------------------------------------------
\begin{comment}
	
	\begin{figure}[H]
\centering
\begin{tikzpicture}[
    % Estilos das camadas
    layer/.style={
        draw, rounded corners=4pt, minimum width=11cm, 
        minimum height=1.4cm, align=center, font=\small\sffamily
    },
    component/.style={
        draw, rounded corners=2pt, minimum width=3.5cm, 
        minimum height=0.9cm, align=center, font=\scriptsize\sffamily
    },
    arrow/.style={-{Stealth[length=3mm]}, thick},
    label text/.style={font=\footnotesize\itshape}
]

% ===== CAMADA 1: APRESENTAÇÃO / ENTRADA =====
\node[layer, fill=blue!15] (entrada) at (0, 6) 
    {Camada de Entrada -- Dados dos Analitos (Cr, Fe, Mn, Ni, Mo)};

% ===== CAMADA 2: PROCESSAMENTO (Java + Python) =====
\node[layer, fill=green!15, minimum height=3.2cm] (processamento) at (0, 3) {};

% Subcamada Java
\node[component, fill=orange!20] (java) at (-3.5, 3.8) 
    {\textbf{Motor Java}\\(Weka / MOA)};
\node[component, fill=orange!10, font=\tiny\sffamily] (tecnicas_j) at (-3.5, 2.8) 
    {RLM $\cdot$ Bagging $\cdot$ Subagging};

% Subcamada Python
\node[component, fill=purple!20] (python) at (3.5, 3.8) 
    {\textbf{Motor Python}\\(scikit-learn)};
\node[component, fill=purple!10, font=\tiny\sffamily] (tecnicas_p) at (3.5, 2.8) 
    {RLM $\cdot$ Bagging $\cdot$ Subagging};

% ===== CAMADA 3: VALIDAÇÃO / PARIDADE =====
\node[layer, fill=yellow!20] (validacao) at (0, 0.5) 
    {Camada de Validação -- Paridade ($\Delta = |Y_{Python} - Y_{Java}| = 0$)};

% ===== SETAS DE FLUXO =====
\draw[arrow] (entrada.south) -- node[right, label text]{CSV} (processamento.north);
\draw[arrow] (processamento.south) -- node[right, label text]{Predições} (validacao.north);

% Conexões internas
\draw[-{Stealth[length=2mm]}, dashed, gray] (java.south) -- (tecnicas_j.north);
\draw[-{Stealth[length=2mm]}, dashed, gray] (python.south) -- (tecnicas_p.north);

\end{tikzpicture}
\caption{Diagrama de arquitetura em camadas do sistema de predição (Java + Python) para os analitos Cr, Fe, Mn, Ni e Mo.}
\label{fig:arquitetura_camadas}
\end{figure}

\end{comment}
%---------------------------------------------------

